\documentclass[11pt]{article}
\usepackage[margin=1.05in]{geometry}
\usepackage{amsmath,amssymb,amsthm,mathtools}
\usepackage{hyperref}
\usepackage{booktabs}
\hypersetup{colorlinks=true,linkcolor=blue,citecolor=blue,urlcolor=blue}

\newtheorem{thm}{Theorem}
\newtheorem{lem}[thm]{Lemma}
\newtheorem{prop}[thm]{Proposition}
\newtheorem{cor}[thm]{Corollary}
\theoremstyle{definition}
\newtheorem{defn}[thm]{Definition}
\newtheorem{rem}[thm]{Remark}
\newtheorem{retr}[thm]{Retraction}

\newcommand{\At}{\widetilde{A}}
\newcommand{\Ah}{\widehat{A}}
\newcommand{\bM}{\mathbf{M}}
\newcommand{\bI}{\mathbf{I}}
\newcommand{\bD}{\mathbf{D}}
\newcommand{\bU}{\mathbf{U}}
\newcommand{\bR}{\mathbf{R}}
\newcommand{\bh}{\mathbf{h}}
\newcommand{\sm}{\smallsetminus}
\newcommand{\EB}{\mathcal{E}^{B}}
\newcommand{\EC}{\mathcal{E}^{C}}
\newcommand{\EX}{\mathcal{E}^{X}}

\title{The gauge theorem in type $B$:\\
one collapse law for every root of $B_N$,\\
and a domain cut out by the root system itself}
\author{Clio}
\date{19 October 2026}

\makeatletter
%% protocol-header.tex -- PROTOCOL §2.3 page-1 block, as a TEMPLATE not an instance.
%%
%% PROTOCOL §2.3: "Every PDF must state, on its first page: author, date,
%% recipient, a title, and the commit hash of the work-in-progress push it
%% corresponds to.  A PDF with no commit hash cannot be checked against the
%% record and should be treated as a draft opinion, not a result."
%%
%% Written 2026-10-16 because the fault had recurred twice and I had twice
%% fixed only the instance:
%%   - 2026-10-14: the Q473 PDF went to Rick with no recipient and no hash.
%%   - 2026-10-14: the Q460 PDF went to Rick saying "Recipient: Robin Langer".
%%     He noticed, and took the commit hash as authority instead.
%% Both times I repaired the one file and wrote "fix the template, not just the
%% instance" in my notes.  This file is that fix.  \input it; do not retype it.
%%
%% USAGE, immediately after \maketitle:
%%   %% protocol-header.tex -- PROTOCOL §2.3 page-1 block, as a TEMPLATE not an instance.
%%
%% PROTOCOL §2.3: "Every PDF must state, on its first page: author, date,
%% recipient, a title, and the commit hash of the work-in-progress push it
%% corresponds to.  A PDF with no commit hash cannot be checked against the
%% record and should be treated as a draft opinion, not a result."
%%
%% Written 2026-10-16 because the fault had recurred twice and I had twice
%% fixed only the instance:
%%   - 2026-10-14: the Q473 PDF went to Rick with no recipient and no hash.
%%   - 2026-10-14: the Q460 PDF went to Rick saying "Recipient: Robin Langer".
%%     He noticed, and took the commit hash as authority instead.
%% Both times I repaired the one file and wrote "fix the template, not just the
%% instance" in my notes.  This file is that fix.  \input it; do not retype it.
%%
%% USAGE, immediately after \maketitle:
%%   %% protocol-header.tex -- PROTOCOL §2.3 page-1 block, as a TEMPLATE not an instance.
%%
%% PROTOCOL §2.3: "Every PDF must state, on its first page: author, date,
%% recipient, a title, and the commit hash of the work-in-progress push it
%% corresponds to.  A PDF with no commit hash cannot be checked against the
%% record and should be treated as a draft opinion, not a result."
%%
%% Written 2026-10-16 because the fault had recurred twice and I had twice
%% fixed only the instance:
%%   - 2026-10-14: the Q473 PDF went to Rick with no recipient and no hash.
%%   - 2026-10-14: the Q460 PDF went to Rick saying "Recipient: Robin Langer".
%%     He noticed, and took the commit hash as authority instead.
%% Both times I repaired the one file and wrote "fix the template, not just the
%% instance" in my notes.  This file is that fix.  \input it; do not retype it.
%%
%% USAGE, immediately after \maketitle:
%%   \input{protocol-header}
%%   \protocolheader{Rick (grandpa-rick)}{clio-vega/proofs}{60e4ce3}
%%
%% The recipient argument is MANDATORY and has no default on purpose: a default
%% is what produced "Recipient: Robin Langer" on a PDF addressed to Rick.

%% \maketitle CLEARS \@date (and \@author, \@title) in the article class, so a
%% macro called AFTER \maketitle that reads \@date prints an EMPTY row.  Caught
%% 2026-10-16 only because I ran pdftotext on page 1 instead of trusting rc=0:
%% the build was clean and the Date field was blank.  Hence: snapshot the date
%% HERE, at \input time in the preamble, where \@date is still live.
\let\protocol@date\@date

\newcommand{\protocolheader}[3]{%
  \begin{center}
  \fbox{\begin{minipage}{0.92\textwidth}
  \small
  \textbf{PROTOCOL §2.3 record.}\\[1mm]
  \begin{tabular}{@{}ll@{}}
  Author     & Clio \\
  Date       & \protocol@date \\
  Recipient  & #1 \\
  Repository & \texttt{#2} \\
  Commit     & \texttt{#3} \\
  \end{tabular}\\[1mm]
  This PDF corresponds to the pushed commit above.  Where this text and the
  commit disagree, \textbf{the commit is authoritative}.
  \end{minipage}}
  \end{center}
  \vspace{2mm}%
}

%%   \protocolheader{Rick (grandpa-rick)}{clio-vega/proofs}{60e4ce3}
%%
%% The recipient argument is MANDATORY and has no default on purpose: a default
%% is what produced "Recipient: Robin Langer" on a PDF addressed to Rick.

%% \maketitle CLEARS \@date (and \@author, \@title) in the article class, so a
%% macro called AFTER \maketitle that reads \@date prints an EMPTY row.  Caught
%% 2026-10-16 only because I ran pdftotext on page 1 instead of trusting rc=0:
%% the build was clean and the Date field was blank.  Hence: snapshot the date
%% HERE, at \input time in the preamble, where \@date is still live.
\let\protocol@date\@date

\newcommand{\protocolheader}[3]{%
  \begin{center}
  \fbox{\begin{minipage}{0.92\textwidth}
  \small
  \textbf{PROTOCOL §2.3 record.}\\[1mm]
  \begin{tabular}{@{}ll@{}}
  Author     & Clio \\
  Date       & \protocol@date \\
  Recipient  & #1 \\
  Repository & \texttt{#2} \\
  Commit     & \texttt{#3} \\
  \end{tabular}\\[1mm]
  This PDF corresponds to the pushed commit above.  Where this text and the
  commit disagree, \textbf{the commit is authoritative}.
  \end{minipage}}
  \end{center}
  \vspace{2mm}%
}

%%   \protocolheader{Rick (grandpa-rick)}{clio-vega/proofs}{60e4ce3}
%%
%% The recipient argument is MANDATORY and has no default on purpose: a default
%% is what produced "Recipient: Robin Langer" on a PDF addressed to Rick.

%% \maketitle CLEARS \@date (and \@author, \@title) in the article class, so a
%% macro called AFTER \maketitle that reads \@date prints an EMPTY row.  Caught
%% 2026-10-16 only because I ran pdftotext on page 1 instead of trusting rc=0:
%% the build was clean and the Date field was blank.  Hence: snapshot the date
%% HERE, at \input time in the preamble, where \@date is still live.
\let\protocol@date\@date

\newcommand{\protocolheader}[3]{%
  \begin{center}
  \fbox{\begin{minipage}{0.92\textwidth}
  \small
  \textbf{PROTOCOL §2.3 record.}\\[1mm]
  \begin{tabular}{@{}ll@{}}
  Author     & Clio \\
  Date       & \protocol@date \\
  Recipient  & #1 \\
  Repository & \texttt{#2} \\
  Commit     & \texttt{#3} \\
  \end{tabular}\\[1mm]
  This PDF corresponds to the pushed commit above.  Where this text and the
  commit disagree, \textbf{the commit is authoritative}.
  \end{minipage}}
  \end{center}
  \vspace{2mm}%
}

\makeatother

\begin{document}
\maketitle
\protocolheader{Rick (grandpa-rick)}{clio-vega/proofs}{124b472}

\begin{abstract}
Let $\psi_X(v)=(1-qv^{h})/(1-q)$ with $h=2$ for $X=C$ and $h=1$ for $X=B$,
$\varphi_X=\psi_X^{-1}$, and $\Ah^{X}_{\bM,\bI}:=\At^{X}_{\bM,\bI}
\prod_{m\in\bM}\psi_X(w_m)$, where $\At^{X}=W^{X}/\EX$ is Motegi's normalised
type-$X$ weight function on the $2^N$ fixed points of $LG(N,2N)$.  This paper
closes gap \textbf{G2} of \texttt{2026-10-18-c2}: for every positive root
$\alpha$ of $B_N$, every $N$, every $\bM$ in the colliding stratum and every
$\bI$, the collapse multiplier is
$\rho^{\alpha}=\prod_{m\in D(\alpha)}\psi_B(w_m)^{-1}$; hence $\Ah^{B}$
collapses with multiplier exactly $1$ at every root of $B_N$.  The
pre-registration was right: the $\iota$-orbit classification is
type-independent, only $h$ moves, and the multiplier of $\Ah$ is again $1$.
Measured $104/104$ at all nine positive roots of $B_2$ and $B_3$ by two
mechanisms, with $124/124$ vacuous strata accounted separately and three
admissible plants firing $13/13$ each.

The proof is one computation, and it puts $h$ in two places rather than one.
\textbf{(1)} Cancelling $W^{B}$'s prefactor $\prod_k(1+u_k^{-1})^{-1}$ against
line~3c's $(u_j+1)$ --- which must be done \emph{analytically}, since the first
has a pole exactly where the second has a zero --- turns the type-$B$ summand
into the type-$C$ summand times $\prod_{\text{heads}}v_j/(v_j+1)$.
Equivalently: \textbf{the type-$B$ head factor is $(v_j-1)$ where type $C$'s is
$(v_j^{2}-1)$} --- that is $1-e^{-\alpha}$ for the boundary root, short in $B$
and long in $C$.  Since the head factor is the same on both sides of an
adjunction ratio, the \emph{unnormalised} ratio is type-independent and Steps
1--3 are not analogous to type $C$'s but identical to them.  \textbf{(2)} The
whole remaining difference is the $u$-diagonal block of the normaliser, and
there the surprise: the factor $\prod_y(1-qy^{2})^{-1}$ produced by Lemma~I2 ---
the lemma that in type $C$ \emph{is} where $\psi$ comes from --- is in type $B$
cancelled exactly by $\EB$'s own $\prod_j(1-qu_j^{2})$, and $\psi_B$ comes
instead from $\EB$'s numerator $\prod_j(1-qu_j)$.  \textbf{The two $\psi$'s have
different origins}; they are not one formula with a parameter.

A pre-registered claim of mine was then \textbf{refuted, and the refutation is
the best thing here.}  I expected the arbitrary-$\iota$-closed theorem to have
content at the second $\tau$-fixed point $\varepsilon=-1$, with
$\varphi_B(-1)=(1-q)/(1+q)$ as a sharp type-discriminator.  It does not:
$(v_j-1)$ vanishes only at $v_j=+1$, so \textbf{type-$B$ head exclusion fails at
$-1$} (measured: head term $\ne0$, $9/9$), and the law fails there --- $0/9$ at
$r\ge1$, against $6/6$ at $r=0$ where there is no head slot at all.  So the
hypothesis is \emph{necessary}, and what replaces my claim is better: the
$\tau$-fixed points at which type-$X$ head exclusion fires are exactly the
$h_X$-th roots of unity, i.e.\ exactly the points of the boundary root
hyperplane $\{w_m^{h}=1\}$.  \textbf{The theorem's domain in each type is cut out
by that type's own root system}, and reaches no further.

Two provenance corrections, both against me.  \textbf{(i)}~Task~0 of my brief
ordered a literature check before proving.  Motegi's
\texttt{normalizedsusythree} (L2767--2771) is the type-$B$ pair multiplier and
\texttt{stabilityWB} (L2779--2783) the type-$B$ boundary one; worse,
\texttt{normalizedsusytwo} (L2618--2622) is the \emph{type-$C$} pair multiplier
--- the headline number of \texttt{2026-10-18-c2}, whose priority section
credited Rick and an earlier session of mine and did not cite Motegi, and which
sits $27$ lines above a statement my own engine docstring already cited.
\textbf{(ii)}~The vacuous strata $r>|s_\alpha\bM|$ vanish by head exclusion, not
by triangularity; \texttt{2026-10-18-c2}'s \texttt{cor:heads} cites (M3) there
and that citation is wrong, though its conclusion is right.
\end{abstract}

\tableofcontents

%=====================================================================
\section{What this session was asked to do, and what happened}

The brief named one theorem, G2, and put a mandatory half-hour of literature
checking in front of it, because the previous session had spent part of itself
proving a statement (G3, $\At\in R$) that was already in Motegi four lines from a
citation I had been handed.  The order was: read the relevant propositions
\emph{first}, and if they force type $B$, say so and grade G2 a corollary rather
than a theorem.  Outcome, in the order it happened:

\begin{enumerate}
\item \textbf{The disambiguation the brief demanded, resolved by the defining
  sentence rather than the phrase.}  Two things in my notes are called ``the
  type $B$ gap''.  \textbf{Q514} asks whether Motegi's \texttt{normalizationWB}
  (L1386) is a product over the positive roots of $B_N$ whose $A$-type factors
  are the same $(1-w_a/w_{i_l})$: it is about the \emph{diagonal}
  $\At_{\bI,\bI}$ and the $A$-type-divisor colength, and it was closed on 10-18.
  \textbf{G2} is about $\rho^{\alpha}$, the \emph{collapse multiplier at a root
  hyperplane}, for all three families.  Different objects; Q514 closed, G2 open.
\item \textbf{G2's conclusion is partly in Motegi, and so is its type-$C$
  twin's.}  \S\ref{sec:priority}.  A real finding, and against me.
\item \textbf{Steps 1--3 do not ``transfer''; they are the same steps}, for a
  reason worth isolating (Lemma~\ref{lem:typecompare}).  \S\ref{sec:proof}.
\item \textbf{The pre-registration on $\rho^{\alpha}$ holds} and the named risk
  --- long and short exchanging between the types --- is explained away rather
  than merely avoided.  \S\ref{sec:thm}, \S\ref{sec:vocab}.
\item \textbf{A pre-registered claim of mine about $\varepsilon=-1$ was
  refuted}, and the replacement is sharper than the claim.
  \S\ref{sec:beyond}.
\item \textbf{One unexpected structural fact}: $\psi_B$ and $\psi_C$ come from
  different blocks.  \S\ref{sec:origin}.
\end{enumerate}

%=====================================================================
\section{Setup}
\label{sec:setup}

Notation follows \texttt{2026-10-18-c2-gauge-theorem-and-q-divisors.tex},
including the corrected divisor nomenclature: divisors are named by the
\emph{root} that cuts them, never by a root length.

Let $R:=\mathbb{C}(q)[w_1^{\pm},\dots,w_N^{\pm}]$, $F:=\operatorname{Frac}R$,
$P:=2^{[N]}$ identified with the $T$-fixed points of $LG(N,2N)$.  Motegi's
symmetrisation formula (\texttt{2610.15344}, eq.\ \texttt{Motherfunction},
L473--480) is
\begin{equation}\label{eq:mother}
   F_{N,n,r}(u\mid w\mid\bI)
   =\frac{1}{(n-r)!}\Bigl(\prod_{k=1}^{n}u_k^{-(N+n)}\Bigr)
     \sum_{\sigma\in S_n}T\bigl(u_{\sigma(1)},\dots,u_{\sigma(n)}\bigr),
\end{equation}
where, writing $v_j:=u_{\sigma(j)}$ and calling slots $1,\dots,r$ \textbf{heads}
and $r+1,\dots,n$ \textbf{tails},
\begin{equation}\label{eq:T}
\begin{aligned}
T(v)=\;&\underbrace{\prod_{j=1}^{r}\Bigl[(1-q)(v_j^2-1)
   \prod_{k=i_j+1}^{N}(v_j-qw_k)\prod_{k=1}^{i_j-1}(v_j-w_k)\Bigr]}_{\text{line 1}}
   \cdot\underbrace{\prod_{1\le j<k\le r}\frac{qv_j-v_k}{v_j-v_k}\,(v_jv_k-1)}_{\text{line 2a}}\\[1mm]
   &\cdot\underbrace{\prod_{j>r}\prod_{k=1}^{N}(v_j-qw_k)}_{\text{line 2b}}
   \cdot\underbrace{\prod_{j>r}\prod_{k\le r}\frac{qv_j-v_k}{v_j-v_k}\,(v_jv_k-1)}_{\text{line 3a}}
   \cdot\underbrace{\prod_{r<j<k\le n}(v_jv_k-q)}_{\text{line 3b}}
   \cdot\underbrace{\prod_{j>r}(\beta v_j-\gamma)}_{\text{line 3c}} .
\end{aligned}
\end{equation}
The two specialisations (L494--501) are
\begin{equation}\label{eq:WBWC}
  W^{C}_{N,n,r}:=F\big|_{\beta=1,\gamma=0},
  \qquad
  W^{B}_{N,n,r}:=\Bigl(\prod_{k=1}^{n}\bigl(1+u_k^{-1}\bigr)^{-1}\Bigr)
                 F\big|_{\beta=1,\gamma=-1},
\end{equation}
and the normalisers (L1057--1062, L1361--1367) are
\begin{align}
  \EC_{N,n,r}(u,w)&=
  \frac{\overbrace{\textstyle\prod_{j,k=1}^{n}\bigl(1-q u_j^{-1}u_k\bigr)}^{\textstyle Q_1}
        \cdot\prod_{j=1}^{n}\prod_{k=1}^{N}(1-qu_jw_k)}
       {\underbrace{\textstyle\prod_{1\le j<k\le n}(1-qu_ju_k)}_{\textstyle Q_2}
        \cdot\prod_{1\le j\le k\le N}(1-qw_jw_k)\cdot(1-q)^{\,n-r}},
  \label{eq:EC}\\[2mm]
  \EB_{N,n,r}(u,w)&=
  \frac{Q_1\cdot\prod_{j=1}^{n}\prod_{k=1}^{N}(1-qu_jw_k)
        \cdot\boxed{\textstyle\prod_{j=1}^{n}(1-qu_j)}}
       {Q_2\cdot\boxed{\textstyle\prod_{j=1}^{n}(1-qu_j^{2})}
        \cdot\prod_{1\le j<k\le N}(1-qw_jw_k)\cdot\prod_{k=1}^{N}(1-qw_k)
        \cdot(1-q)^{\,n-r}} .
  \label{eq:EB}
\end{align}
Write $\At^{X}_{N,n,r}:=W^{X}_{N,n,r}/\EX_{N,n,r}$ (Motegi's
$\widetilde{WX}$) and, for a multiset $\bU$ of values,
$\At^{X}(\bU):=\At^{X}_{N,|\bU|,r}(\bU\mid w\mid\bI)$.  All of
\eqref{eq:mother}--\eqref{eq:EB} was re-transcribed first-hand this session from
the e-print source at the line numbers given; the five local copies of that
source agree (md5 \texttt{c15a23f2}, $5/5$).

\begin{rem}[the two boxed blocks]
\label{rem:boxed}
$\EB$ and $\EC$ differ in four places, but two of them --- the $w$-side
$\prod_{j<k}$ versus $\prod_{j\le k}$, and the extra $\prod_k(1-qw_k)$ --- are
independent of $n$ and of $u$, hence cancel identically in any ratio
$\EX_{n}/\EX_{s}$ at fixed $w$.  The only surviving difference is the pair of
boxed $u$-diagonal blocks.  Measured, symbolically and for all $\bM$ at
$N=2,3$: $\EB/\EC=\prod_{k\notin\bM}(1-qw_k^{2})/(1-qw_k)$ under
$u=w_{\bM}$, $28/28$.
\end{rem}

For $\bM\in P$ put $\At^{X}_{\bM,\bI}:=\At^{X}(w_{\bM})$ with
$w_{\bM}=\{w_m:m\in\bM\}$ and $n=s:=|\bM|$; the absence of trailing $1$'s is
legitimate in type $B$ by \emph{Motegi's own type-$B$} iterated stability,
L3602--3608 (for $s\ge r$), and $\At^{B}_{\bM,\bI}=0$ for $s<r$ is L3620--3626.
I flag this because the corresponding type-$C$ citations, L2822--2828 and L2843,
are the ones carried in my engine's docstring, and reusing a type-$C$ locator
for a type-$B$ claim is exactly the 2026-10-17 defect.

\begin{defn}[the collision data of a root of $B_N$]\label{def:D}
The positive roots of $B_N$ are $e_i-e_j$, $e_i+e_j$ $(i<j)$ and $e_m$.  For
each, $H_\alpha$ is the vanishing locus of the corresponding $T$-curve
character, $s_\alpha$ the signed-permutation action on $P$, and $D(\alpha)$ the
set of labels deleted:
\[
\begin{array}{llll}
\text{root} & H_\alpha & s_\alpha(\bM)\ \text{on the colliding stratum} & D(\alpha)\\\hline
e_i-e_j & \{w_i=w_j\} & \bM\triangle\{i,j\}\quad(|\bM\cap\{i,j\}|=1) & \varnothing\\
e_m  & \{w_m=1\}&\bM\sm\{m\}\quad(m\in\bM) & \{m\}\\
e_a+e_b & \{w_aw_b=1\}&\bM\sm\{a,b\}\quad(\{a,b\}\subseteq\bM) & \{a,b\}
\end{array}
\]
\end{defn}

\noindent
This is Definition~\texttt{def:D} of \texttt{2026-10-18-c2} at $h=1$: the
boundary root of $C_N$ is $2e_m$ with character $w_m^{2}$, of $B_N$ is $e_m$ with
character $w_m$.  Lemma~\texttt{lem:model} there --- that this is the
signed-permutation action --- is a statement about $W(C_N)=W(B_N)$ and is
type-independent.

Set $\iota(v)=v^{-1}$.  A finite multiset $\bD\subset F^{\times}$ is
\textbf{$\iota$-closed} if some involution $\tau$ of its index set satisfies
$y_{\tau(j)}=y_j^{-1}$; the $\tau$-fixed points are the $y$ with $y^{2}=1$.

\begin{lem}[the three root families of $B_N$ are the three $\iota$-orbit types]
\label{lem:orbits}
Let $\alpha>0$ in $B_N$, $\bM$ in its colliding stratum, and
$\bD(\alpha):=\{w_m:m\in D(\alpha)\}$ restricted to $H_\alpha$.  Then
$\bD(\alpha)$ is $\iota$-closed, and
\[
  \alpha=e_i-e_j\ \Rightarrow\ \bD=\varnothing;\quad
  \alpha=e_m\ \Rightarrow\ \bD=\{1\}\ (\tau\text{-fixed});\quad
  \alpha=e_a+e_b\ \Rightarrow\ \bD=\{x,x^{-1}\}\ (\text{free }\tau\text{-orbit}),
\]
and on $H_\alpha$ the value multisets satisfy
$w_{\bM}=w_{s_\alpha\bM}\sqcup\bD(\alpha)$.
\end{lem}

\begin{proof}
Verbatim the proof of Lemma~\texttt{lem:orbits} of \texttt{2026-10-18-c2}, with
one simplification: where type $C$'s boundary locus $\{w_m^{2}=1\}$ gives
$\varepsilon=w_m$ with $\varepsilon=\pm1$, type $B$'s locus $\{w_m=1\}$ gives
$\varepsilon=1$ outright.  The $A$-type clause is that on $\{w_i=w_j\}$ with
$i\in\bM\not\ni j$ the two value multisets are \emph{equal}.
\end{proof}

\begin{rem}[the one structural asymmetry between the types]
\label{rem:components}
$\{w_m^{2}=1\}$ has the two components $w_m=\pm1$; $\{w_m=1\}$ has one.  Both
$+1$ and $-1$ are $\tau$-fixed points of $\iota$ in either type, but only in
type $C$ are both of them \emph{root} hyperplanes.  \S\ref{sec:beyond} shows
that this is not a cosmetic difference: it is exactly where the type-$B$
theorem's domain stops.
\end{rem}

%=====================================================================
\section{Steps 1--3 are not analogous to type $C$'s; they are identical}
\label{sec:proof}

\begin{lem}[the type comparison, with the prefactor cancelled analytically]
\label{lem:typecompare}
Fix $\sigma\in S_n$ and write $T^{X}_\sigma$ for the summand of
\eqref{eq:mother} contributing to $W^{X}$, so that
$W^{X}=\frac{1}{(n-r)!}\bigl(\prod_k u_k^{-(N+n)}\bigr)\sum_\sigma T^{X}_\sigma$
with the convention that $W^{B}$'s prefactor is distributed over the summands.
Then, as rational functions,
\[
  T^{B}_\sigma=T^{C}_\sigma\cdot\prod_{j=1}^{r}\frac{v_j}{v_j+1},
\]
and explicitly $T^{B}_\sigma$ equals \eqref{eq:T} with line~3c deleted, line~1's
head factor $(1-q)(v_j^{2}-1)$ replaced by
\[
  \boxed{\ (1-q)(v_j-1)\ }
\]
and an overall $\prod_{k=1}^{n}u_k$.  In particular $T^{B}_\sigma$ has no pole at
$u_k=-1$.
\end{lem}

\begin{proof}
$\prod_{k=1}^n(1+u_k^{-1})^{-1}=\prod_{k=1}^n u_k/(u_k+1)$, and line~3c at
$\beta=1,\gamma=-1$ is $\prod_{j>r}(v_j+1)$.  Since
$\prod_{k=1}^{n}(u_k+1)=\prod_{j=1}^{n}(v_j+1)$ --- the same multiset of values
--- we get
\[
  \prod_{k=1}^{n}\frac{u_k}{u_k+1}\cdot\prod_{j>r}(v_j+1)
  =\frac{\prod_{k=1}^{n}u_k}{\prod_{j\le r}(v_j+1)} ,
\]
so $T^{B}_\sigma=[\text{lines }1,2a,2b,3a,3b]\cdot\prod_k u_k\big/
\prod_{j\le r}(v_j+1)$.  Absorbing $1/(v_j+1)$ into line~1 via
$v_j^{2}-1=(v_j-1)(v_j+1)$ gives the boxed head factor and the displayed
explicit form.  Comparing with
$T^{C}_\sigma=[\text{lines }1,2a,2b,3a,3b]\cdot\prod_{j>r}v_j$ and using
$\prod_k u_k=\prod_{j\le r}v_j\cdot\prod_{j>r}v_j$ gives the ratio.  No factor of
the explicit form has $u_k+1$ in a denominator, whence the last claim.
\end{proof}

\begin{rem}[why this cancellation has to be done before substituting]
\label{rem:cancelfirst}
$W^{B}$'s prefactor has a \emph{pole} at $u_k=-1$ and line~3c a \emph{zero}
there, so the literal expression \eqref{eq:WBWC} evaluates to $\infty\cdot0$ at
$u_k=-1$: undefined, not wrong.  I ran the type-independence census on the
literal form first and it returned $0/18$ at $\varepsilon=-1$, which I nearly
read as a refutation of Lemma~\ref{lem:typeindep}.  It was an undefined
computation.  The correct reading of $0/18$ was ``you substituted into a
removable singularity''; and when I then did the cancellation, the real
phenomenon turned out to be \S\ref{sec:beyond}, which is a different and
sharper statement.  Both the false reading and the true one live at the same
point of the same locus.
\end{rem}

\begin{lem}[head exclusion in type $B$, and exactly where it fails]
\label{lem:head}
Let $\bD$ be $\iota$-closed with realising involution $\tau$, and let the values
of $\bR\sqcup\bD$ be \textbf{pairwise distinct}.  Let $\sigma$ place slot $p$ of
$\bD$ in a head position.  Then $T^{B}_\sigma=0$ if either
\begin{enumerate}
\item[(a)] $\tau(p)\ne p$, or
\item[(b)] $\tau(p)=p$ and $u_p=+1$;
\end{enumerate}
and if $\tau(p)=p$ with $u_p=-1$ then $T^{B}_\sigma\ne0$ in general.  In type
$C$, both (a) and $\tau(p)=p$ with $u_p=\pm1$ give $T^{C}_\sigma=0$.
\end{lem}

\begin{proof}
(a) The slot pair $\{p,\tau(p)\}$ contains the head $p$, so it carries the
factor $(u_pu_{\tau(p)}-1)$ from line~2a (if $\tau(p)$ is also a head) or
line~3a (if it is a tail); and $u_pu_{\tau(p)}=1$.  These two lines are
unchanged by Lemma~\ref{lem:typecompare}, so this case is type-independent.
Distinctness enters exactly here: the same slot pair carries the
\emph{denominator} $(v_p-v_{p'})$ in lines 2a/3a, and if two values of $\bD$
coincide the term is $0/0$ rather than $0$
(Remark~\texttt{rem:distinct} of \texttt{2026-10-18-c2}).

(b) and the failure.  By Lemma~\ref{lem:typecompare} the head factor at slot $p$
is $(1-q)(u_p-1)$ in type $B$ and $(1-q)(u_p^{2}-1)$ in type $C$.  A
$\tau$-fixed point has $u_p^{2}=1$, i.e.\ $u_p=\pm1$.  Then $u_p^{2}-1=0$
always, while $u_p-1=0$ iff $u_p=+1$ and equals $-2$ at $u_p=-1$; and no other
factor of the explicit form of Lemma~\ref{lem:typecompare} vanishes at $u_p=-1$
for generic $\bR$ and $w$.
\end{proof}

\begin{rem}[measured]
Lemma~\ref{lem:head}'s split is exactly what the head-slot term does: at
$\varepsilon=+1$ it vanishes in $9/9$ instances ($N=2,3$, $r=1,2$, every $\bI$),
and at $\varepsilon=-1$ it fails to vanish in $9/9$.  A clean split on the
nose of the stated discriminator.
\end{rem}

\begin{defn}[admissible]\label{def:adm}
Call an $\iota$-closed $\bD$ \textbf{$X$-admissible} if its values are pairwise
distinct and, when $X=B$, no $\tau$-fixed point of $\bD$ equals $-1$.  (For
$X=C$ no extra condition.)  When $r=0$ every $\iota$-closed $\bD$ with distinct
values is admissible in both types, since there is then no head slot.
\end{defn}

\begin{cor}[the head selections match]\label{cor:heads}
Let $\bD$ be $X$-admissible, $s=|\bR|$, $d=|\bD|$, $n=s+d$.  Then
\[
  \frac{1}{(n-r)!}\sum_{\sigma\in S_n}T^{X}_\sigma
  =\sum_{\substack{(p_1,\dots,p_r)\ \textup{distinct}\\ \textup{slots of }\bR}}
      T^{X}\bigl(p_1,\dots,p_r;\ \textup{tails}=\textup{the rest}\bigr),
\]
and the index set on the right is the same for $\bR\sqcup\bD$ and for $\bR$.
\end{cor}

\begin{proof}
Tail-symmetry: lines 2b, 3b are tail-symmetric, line 3a is a product over tails
of head-dependent factors, and the extra data of Lemma~\ref{lem:typecompare}
--- the overall $\prod_k u_k$ and the head factors --- is symmetric in the tails.
So the $(n-r)!$ tail orderings give equal terms.  By Lemma~\ref{lem:head}, with
$\bD$ admissible, the surviving head selections avoid $\bD$.
\end{proof}

\begin{lem}[the unnormalised adjunction ratio is type-independent]
\label{lem:typeindep}
Let $\bD$ be $B$- and $C$-admissible and $s=|\bR|\ge r$.  Then
\[
  \frac{W^{B}_{N,n,r}(\bR\sqcup\bD)}{W^{B}_{N,s,r}(\bR)}
  =\frac{W^{C}_{N,n,r}(\bR\sqcup\bD)}{W^{C}_{N,s,r}(\bR)} .
\]
\end{lem}

\begin{proof}
By Corollary~\ref{cor:heads} both numerator and denominator are sums over the
\emph{same} head selections $\bh\subset\bR$.  Fix one.  By
Lemma~\ref{lem:typecompare} the type-$B$ summand is the type-$C$ summand times
$\prod_{v\in\bh}v/(v+1)$, a factor depending only on $\bh$ --- and in particular
not on $\bD$.  It therefore cancels between numerator and denominator term by
term.
\end{proof}

\begin{rem}
This is stronger than yesterday's gap statement, which read ``Steps 1 and 2 are
untouched; what changes is the bookkeeping of Steps 3--5''.  Steps 1, 2
\emph{and} 3 are untouched: the whole $F$-side computation is the type-$C$ one,
and only the normaliser has to be redone.  I had located the type-dependence one
step too early --- and, per Lemma~\ref{lem:head}, one \emph{hypothesis} too
weakly.
\end{rem}

\noindent
So I may quote eq.~\texttt{eq:Fratio} of \texttt{2026-10-18-c2} as it stands.
Recall
\[
  Q_1(\bD)=\!\!\prod_{\text{ordered }(j,k)}\!\!\Bigl(1-q\tfrac{y_k}{y_j}\Bigr),\quad
  Q_2(\bD)=\!\!\prod_{\{j,k\},\,j\ne k}\!\!(1-qy_jy_k),\quad
  P_2(\bD)=\!\!\prod_{\{j,k\},\,j\ne k}\!\!(y_jy_k-q),
\]
and the two reindexing identities, proved there for arbitrary $\iota$-closed
$\bD$ and \emph{type-free} (they are statements about a multiset with an
involution; no weight function occurs in either):
\begin{equation}\label{eq:I1I2}
  \text{(I1)}\quad P_2(\bD)=\Bigl(\prod_{y\in\bD}y\Bigr)^{d-1}Q_2(\bD),
  \qquad
  \text{(I2)}\quad \frac{Q_2(\bD)^{2}}{Q_1(\bD)}=\frac{1}{\prod_{y\in\bD}(1-qy^{2})} .
\end{equation}

\begin{prop}[Step 3, quoted]\label{prop:step3}
With $\bD$ admissible, $s\ge r$, $n=s+d$, and either type,
\begin{equation}\label{eq:Fratio}
  \frac{W^{X}_{N,n,r}(\bR\sqcup\bD)}{W^{X}_{N,s,r}(\bR)}
  =\Bigl(\prod_{y\in\bD}y\Bigr)^{1-d} P_2(\bD)
   \prod_{y\in\bD}\Bigl[\prod_{v\in\bR}\bigl(1-\tfrac{qy}{v}\bigr)
                        \prod_{k=1}^{N}(1-qyw_k)\Bigr].
\end{equation}
\end{prop}

\begin{proof}
For $X=C$ this is Steps 1--3 of Theorem~\texttt{thm:gauge} in
\texttt{2026-10-18-c2}; for $X=B$ it follows by Lemma~\ref{lem:typeindep}.
\end{proof}

%=====================================================================
\section{Step 4: the normaliser, and the theorem}
\label{sec:thm}

\begin{lem}[Step 4, type $B$]\label{lem:step4}
With $\bD$ admissible and $s\ge r$,
\begin{equation}\label{eq:EBratio}
  \frac{\EB_{N,n,r}}{\EB_{N,s,r}}
  =\frac{Q_1(\bD)}{Q_2(\bD)\,(1-q)^{d}}
   \cdot\frac{\prod_{y\in\bD}(1-qy)}{\prod_{y\in\bD}(1-qy^{2})}
   \cdot\prod_{y\in\bD}\Bigl[\prod_{v\in\bR}\bigl(1-\tfrac{qy}{v}\bigr)
                        \prod_{k=1}^{N}(1-qyw_k)\Bigr].
\end{equation}
\end{lem}

\begin{proof}
Both normalisers sit at the same alphabet $w$, so by Remark~\ref{rem:boxed} the
two $w$-only blocks cancel.  Adjoining $\bD$ changes each remaining block thus.

$Q_1$: acquires the $\bD\bD$ pairs $Q_1(\bD)$ and the mixed pairs
$\prod_{y\in\bD}\prod_{v\in\bR}\bigl(1-q\tfrac{v}{y}\bigr)
 \bigl(1-q\tfrac{y}{v}\bigr)$, numerator.

$\prod_j\prod_k(1-qu_jw_k)$: acquires $\prod_{y}\prod_{k}(1-qyw_k)$, numerator.

$Q_2$: acquires $Q_2(\bD)\prod_{y}\prod_{v\in\bR}(1-qyv)$, denominator.

$(1-q)^{\,n-r}$: acquires $(1-q)^{d}$, denominator.

\emph{The boxed blocks:} $\prod_j(1-qu_j)$ acquires $\prod_{y}(1-qy)$ in the
numerator, and $\prod_j(1-qu_j^{2})$ acquires $\prod_{y}(1-qy^{2})$ in the
denominator.

Finally, since $\tau$ is an involution with $y_{\tau(j)}=y_j^{-1}$,
$\prod_{y}\prod_{v}\bigl(1-q\tfrac{v}{y}\bigr)=\prod_{y}\prod_{v}(1-qvy)$, which
cancels the denominator's $\prod_{y}\prod_{v}(1-qyv)$.  Collecting gives
\eqref{eq:EBratio}.
\end{proof}

\begin{thm}[the type-$B$ $\iota$-adjunction law]\label{thm:gaugeB}
Let $\bR$ be a finite multiset of values with $s=|\bR|\ge r=|\bI|$ and let
$\bD$ be $\iota$-closed, $d=|\bD|$, and \textbf{$B$-admissible}
(Definition~\ref{def:adm}): values of $\bR\sqcup\bD$ pairwise distinct, and no
$\tau$-fixed point of $\bD$ equal to $-1$ unless $r=0$.  Then, at a generic
point of the locus on which $\bD$ is $\iota$-closed,
\[
  \boxed{\ \At^{B}(\bR\sqcup\bD)=\Bigl(\prod_{y\in\bD}\varphi_B(y)\Bigr)\At^{B}(\bR),
  \qquad \varphi_B(y)=\frac{1-q}{1-qy}\ }
\]
\end{thm}

\begin{proof}
Divide \eqref{eq:Fratio} by \eqref{eq:EBratio}; the bracketed products are
identical and cancel, leaving
\[
  \frac{\At^{B}(\bR\sqcup\bD)}{\At^{B}(\bR)}
  =\Bigl(\prod_{y}y\Bigr)^{1-d}P_2(\bD)\cdot
   \frac{Q_2(\bD)\,(1-q)^{d}}{Q_1(\bD)}\cdot
   \frac{\prod_{y}(1-qy^{2})}{\prod_{y}(1-qy)} .
\]
By (I1), $\bigl(\prod_y y\bigr)^{1-d}P_2(\bD)=Q_2(\bD)$, so the first three
factors are $(1-q)^{d}Q_2(\bD)^{2}/Q_1(\bD)$, which by (I2) is
$(1-q)^{d}\big/\prod_{y}(1-qy^{2})$.  Hence
\[
  \frac{\At^{B}(\bR\sqcup\bD)}{\At^{B}(\bR)}
  =\frac{(1-q)^{d}}{\prod_{y}(1-qy^{2})}\cdot
   \frac{\prod_{y}(1-qy^{2})}{\prod_{y}(1-qy)}
  =\prod_{y\in\bD}\frac{1-q}{1-qy}. \qedhere
\]
\end{proof}

\begin{cor}[G2: $\rho^{\alpha}$ for every root of $B_N$, all $N$]\label{cor:rhoB}
Set $\psi_B(v)=(1-qv)/(1-q)$, $\Psi_B(\bM)=\prod_{m\in\bM}\psi_B(w_m)$ and
$\Ah^{B}_{\bM,\bI}=\At^{B}_{\bM,\bI}\Psi_B(\bM)$.  For every $N$, every positive
root $\alpha$ of $B_N$, every $\bM$ in the colliding stratum of $\alpha$ and
every $\bI$, on $H_\alpha$
\[
   \At^{B}_{\bM,\bI}=\rho^{\alpha}\,\At^{B}_{s_\alpha\bM,\bI},
   \qquad
   \rho^{\alpha}=\prod_{m\in D(\alpha)}\varphi_B(w_m)
                =\prod_{m\in D(\alpha)}\psi_B(w_m)^{-1},
\]
and consequently $\Ah^{B}_{\bM,\bI}=\Ah^{B}_{s_\alpha\bM,\bI}$ on $H_\alpha$:
the collapse multiplier of $\Ah^{B}$ is exactly $1$ at every root of $B_N$, for
every $N$.
\end{cor}

\begin{proof}
Put $\bR=w_{s_\alpha\bM}$, $\bD=\bD(\alpha)$, so $w_{\bM}=\bR\sqcup\bD$ with
$\bD$ $\iota$-closed (Lemma~\ref{lem:orbits}).  Suppose $|\bR|\ge r$.
\emph{$\bD(\alpha)$ is $B$-admissible for every root of $B_N$}: its values are
$\varnothing$, $\{1\}$ or $\{x,x^{-1}\}$, so its only possible $\tau$-fixed
point is $+1$ --- never $-1$ --- and distinctness holds at the generic point of
$H_\alpha$ (the values of $\bR$ are distinct indeterminates; $|\bD|\le2$; and
when $|\bD|=2$ its values satisfy $w_aw_b=1$ and so coincide only on
$\{w_a^{2}=1\}\cap H_{e_a+e_b}$, a proper closed subvariety).  This is the
point at which Lemma~\ref{lem:head}'s restriction is seen not to bite:
\emph{no root of $B_N$ produces $-1$ as a deleted value}, precisely because
$\{w_m=-1\}$ is not a root hyperplane of $B_N$ (Remark~\ref{rem:components}).
Theorem~\ref{thm:gaugeB} then gives the multiplier at the generic point; both
sides lie in $\mathcal{O}_{H_\alpha}$, so the identity holds in the local ring.
Evaluating $\varphi_B$:
\[
\begin{array}{llll}
\alpha & |D| & \bD\ \text{on}\ H_\alpha & \rho^{\alpha}\\\hline
e_i-e_j & 0 & \varnothing & 1\\
e_m & 1 & \{1\} & \varphi_B(1)=\dfrac{1-q}{1-q}=1\\[2mm]
e_a+e_b & 2 & \{x,x^{-1}\} & \dfrac{(1-q)^{2}}{(1-qx)(1-qx^{-1})}
\end{array}
\]
If $|\bR|<r$ both sides vanish: Proposition~\ref{prop:vac}.  For $\Ah^{B}$:
$\bM$ and $s_\alpha\bM$ differ exactly in $D(\alpha)$, so
$\Psi_B(\bM)=\Psi_B(s_\alpha\bM)\prod_{m\in D(\alpha)}\psi_B(w_m)$ and the two
factors are mutually inverse.  For $\alpha=e_i-e_j$, $D=\varnothing$ and
$\Psi_B(\bM)=\Psi_B(s_{ij}\bM)$ on $\{w_i=w_j\}$ since $\psi_B(w_i)=\psi_B(w_j)$
there.
\end{proof}

\begin{rem}[the $d=0$ case is $u$-symmetry, not an instance of the box]
At $\bD=\varnothing$ Theorem~\ref{thm:gaugeB} reads $\At^{B}(\bR)=\At^{B}(\bR)$,
vacuous as an identity; the $A$-type content sits entirely in the second clause
of Lemma~\ref{lem:orbits}, that the two value multisets coincide on $H_\alpha$.
As in type $C$ I say this rather than claim that all three families are
instances of one identity.
\end{rem}

\begin{prop}[the vacuous strata vanish by head exclusion, not by triangularity]
\label{prop:vac}
Let $\alpha>0$, $\bM$ in its colliding stratum, $\bR=w_{s_\alpha\bM}$ and
$r=|\bI|>|\bR|$.  Then on $H_\alpha$ both $\At^{X}_{s_\alpha\bM,\bI}=0$ and
$\At^{X}_{\bM,\bI}=0$, and Corollary~\ref{cor:rhoB} reads
$0=\rho^{\alpha}\cdot0$.
\end{prop}

\begin{proof}
$\At^{X}_{s_\alpha\bM,\bI}=0$ is lower-cardinality vanishing
($|s_\alpha\bM|<r$): L3620--3626 in type $B$, L2843 in type $C$.  For
$\At^{X}_{\bM,\bI}$: by Lemma~\ref{lem:head} and Corollary~\ref{cor:heads} the
surviving head selections are the $r$-element subsets of the slots of $\bR$, of
which there are none when $r>|\bR|$; the sum \eqref{eq:mother} is empty.
\end{proof}

\begin{retr}[a wrong citation in my own type-$C$ paper]\label{retr:M3}
Corollary~\texttt{cor:heads} of \texttt{2026-10-18-c2} disposes of this case with
\emph{``(If $s<r$ both sides vanish: $\At_{\bM,\bI}=\At_{\bR,\bI}=0$ by
\textup{(M3)}\ldots)''}.  For $\At_{\bR,\bI}$ that is right.  For
$\At_{\bM,\bI}$ it is \textbf{wrong}: $|\bM|=|\bR|+d$, so in the pair family with
$|\bR|=r-1$ one has $|\bM|=r+1>r$ and (M3)'s cardinality clause does not apply,
while its triangularity clause $m_{|\bM|-r+j}\ge i_j$ is generally satisfied ---
e.g.\ $N=2$, $\bM=\{1,2\}$, $\bI=\{1\}$, where $m_2=2\ge1$.  The conclusion is
correct and the reason given is not; the correct reason is head exclusion.  I
single this out because it is the characteristic shape of my errors this month: a
true statement resting on a cited authority that does not cover it, and no
objector, because the conclusion is right.  Measured here: $124/124$ instances
at $N=2,3$ over all nine roots vanish as Proposition~\ref{prop:vac} says.
\end{retr}

%=====================================================================
\section{Where the two $\psi$'s come from, and why that is not one mechanism}
\label{sec:origin}

Yesterday I answered Rick's question \emph{``where does $\psi$ come from on
Motegi's side?''}\ with (I2): $\varphi=\psi^{-1}$ is $(1-q)/(1-qy^{2})$ because
$Q_2^{2}/Q_1=\prod_y(1-qy^{2})^{-1}$, the identity holding because $\tau$ matches
$\EC$'s $\prod_{j,k}(1-qu_j^{-1}u_k)$ against its $\prod_{j<k}(1-qu_ju_k)$; and
I described this as ``$\psi$ is a ratio of two blocks of the normaliser, read
along $\iota$''.  Theorem~\ref{thm:gaugeB} shows that answer is
\textbf{type-$C$-specific}.  Compare the Step-5 computations block by block:
\begin{center}
\begin{tabular}{@{}lll@{}}
\toprule
 & type $C$ & type $B$\\
\midrule
(I1)+(I2) output & $(1-q)^{d}\big/\prod_y(1-qy^{2})$ & $(1-q)^{d}\big/\prod_y(1-qy^{2})$\\
$\EX$'s $u$-diagonal blocks & \emph{none} & $\prod_y(1-qy^{2})\big/\prod_y(1-qy)$\\
\midrule
$\prod_y\varphi_X(y)$ & $(1-q)^{d}\big/\prod_y(1-qy^{2})$ & $(1-q)^{d}\big/\prod_y(1-qy)$\\
\bottomrule
\end{tabular}
\end{center}
In type $B$ the output of (I2) --- the very factor that \emph{is} $\psi_C$ --- is
cancelled \textbf{exactly} by $\EB$'s $\prod_j(1-qu_j^{2})$, and $\psi_B$ is
instead $\EB$'s numerator block $\prod_j(1-qu_j)$, read along $\iota$.

\begin{cor}[the type-dependence of $\rho$, isolated]\label{cor:ratio}
For a positive root $\alpha$ of $B_N$ resp.\ $C_N$ with the same collision data
$D(\alpha)$,
\[
  \frac{\rho^{\alpha}_{B}}{\rho^{\alpha}_{C}}
  =\prod_{m\in D(\alpha)}\frac{1-qw_m^{2}}{1-qw_m},
\]
which is precisely the ratio of the boxed $u$-diagonal blocks of
\eqref{eq:EB} and \eqref{eq:EC}.  More generally, for $\bD$ admissible in both
types, $\prod_{y}\varphi_B(y)\big/\prod_{y}\varphi_C(y)
 =\prod_{y\in\bD}(1-qy^{2})/(1-qy)$.
\end{cor}

\begin{proof}
By Lemma~\ref{lem:typeindep} the $W$-ratios agree, so the quotient of the
multipliers is the inverse quotient of the $\EX$-ratios, which by
Remark~\ref{rem:boxed} and Lemma~\ref{lem:step4} is the boxed-block ratio.
\end{proof}

\begin{rem}[the honest revision of my answer to G5]\label{rem:G5}
The uniform formula $\psi_X(v)=(1-qv^{h})/(1-q)$ is correct and it is a
\emph{coincidence of outputs}, not one derivation with a parameter.  No step of
either proof has $h$ as a variable.  What there is: (a) a type-free identity
(I2) with output $\prod(1-qy^{2})^{-1}$, and (b) the $u$-diagonal blocks the
normaliser happens to carry.  In type $C$ (a) survives and (b) is empty; in type
$B$ (b) annihilates (a) and supplies the answer itself.  So the general
statement of ``where $\psi$ comes from'' is \emph{the $u$-diagonal content of
$\EX$ after the $\iota$-pairing}, which reduces to (I2) only when $\EX$ has no
$u$-diagonal block of its own.  And $h$ turns out to live in \emph{two} places,
not one: also in the head factor $(v^{h}-1)$ of Lemma~\ref{lem:typecompare},
which is what \S\ref{sec:beyond} is about.
\end{rem}

%=====================================================================
\section{A pre-registered claim of mine, refuted; and the domain is the root system}
\label{sec:beyond}

Proposition~\ref{prop:specialise} below shows Motegi's propositions already
cover $\bD$ of the two shapes $\{1\}$ and $\{x,x^{-1}\}$.  I pre-registered,
when drafting this paper, that Theorem~\ref{thm:gaugeB}'s extra content would be
visible at the second $\tau$-fixed point:

\begin{quote}
\emph{(pre-registered, and \textbf{false})} \ $\bD=\{-1\}$ is $\iota$-closed,
so $\At^{B}(\bR\sqcup\{-1\})=\frac{1-q}{1+q}\At^{B}(\bR)$ whereas
$\At^{C}(\bR\sqcup\{-1\})=\At^{C}(\bR)$; this is the sharp type-discriminator,
in neither of Motegi's propositions.
\end{quote}

\begin{thm}[the $B$-admissibility hypothesis is necessary]\label{thm:necessary}
For $\bD=\{-1\}$ and $r\ge1$ the conclusion of Theorem~\ref{thm:gaugeB} is
\textbf{false}.  For $r=0$ it holds.
\end{thm}

\begin{proof}
For $r=0$ there is no head slot, Lemma~\ref{lem:head} is vacuous, and the proof
of Theorem~\ref{thm:gaugeB} goes through: $\bD$ is admissible by
Definition~\ref{def:adm}.  For $r\ge1$, Lemma~\ref{lem:head} shows that the
$\sigma$ placing $-1$ in a head slot contributes a nonzero term, so
Corollary~\ref{cor:heads}'s index sets do \emph{not} match and
Lemma~\ref{lem:typeindep} is unavailable.  That is an obstruction to the proof,
not to the statement; the statement is refuted by computation.  Measured with
the cancelled form of Lemma~\ref{lem:typecompare}, symbolic in $w$ and $q$, over
$N=2,3$, $|\bR|\le2$ and all $\bI$: \textbf{$r=0$: $6/6$ agree with
$\varphi_B(-1)$; $r\ge1$: $0/14$}, and the measured ratio is not even a function
of $q$ alone --- at
$N=2$, $s=1$, $\bI=\{1\}$ it is
\[
  \frac{18q^{3}w_1-45q^{2}w_1-108q^{2}+27qw_1+45q+63}
       {10q^{3}w_1+45q^{2}w_1+10q^{2}+35qw_1+45q+35},
\]
whereas $\varphi_B(-1)=(1-q)/(1+q)$ depends on no $w$ at all.  So the failure is
not a different scalar; the adjunction is not by a scalar.
\end{proof}

\begin{rem}[the shape of this refutation is one I have seen before]
A hypothesis of Lemma~\ref{lem:head} was discovered by its own failure mode, and
\textbf{the discriminator is a head slot --- the discriminator the lemma itself
names}.  That is verbatim the structure of
Remark~\texttt{rem:distinct} of \texttt{2026-10-18-c2}, where the omitted
distinctness hypothesis was refuted by a census that was green at $r=0$ and
failed at $r\ge1$.  Two different hypotheses of the same lemma, found the same
way, one day apart.  The lesson I draw is not about either hypothesis: it is that
\emph{$r=0$ is a degenerate stratum which silently satisfies every head
hypothesis}, so a census that is green at $r=0$ and never run at $r\ge1$ carries
no information about head exclusion at all.
\end{rem}

\noindent
What replaces the pre-registered claim is better, because it says why $-1$ had to
be the place where this happens.

\begin{cor}[the theorem's domain is cut out by its own root system]
\label{cor:domain}
In type $X$ the $\tau$-fixed points at which head exclusion fires are exactly
the $h_X$-th roots of unity: $\{\pm1\}$ for $h=2$ (type $C$), $\{+1\}$ for $h=1$
(type $B$).  These are exactly the points of the boundary root hyperplane
$\{w_m^{h}=1\}$ of that type.  Hence
\begin{enumerate}
\item every root hyperplane of $B_N$ delivers an admissible $\bD$
  (Corollary~\ref{cor:rhoB}), so G2 is untouched by
  Theorem~\ref{thm:necessary};
\item the type-$B$ theorem reaches no $\tau$-fixed point beyond its own root
  system, and the type-$C$ theorem reaches both components of $\{w_m^{2}=1\}$
  --- which \emph{is} content beyond \texttt{stabilityWC}, since that
  proposition is stated at $u_n=1$ only.
\end{enumerate}
\end{cor}

\begin{proof}
A $\tau$-fixed point satisfies $y^{2}=1$.  By Lemma~\ref{lem:typecompare} the
head factor is $(1-q)(v^{h}-1)$ with $h=h_X$ --- $(v^{2}-1)$ in type $C$,
$(v-1)$ in type $B$ --- whose zero set among $\{y:y^{2}=1\}$ is the set of
$h$-th roots of unity; and $\{w_m^{h}=1\}$ is by Definition~\ref{def:D} the
boundary root hyperplane.  (1) is the admissibility computation in
Corollary~\ref{cor:rhoB}.  (2) is Lemma~\ref{lem:head}'s type-$C$ clause.
\end{proof}

\begin{rem}[so what \emph{is} new beyond Motegi's propositions]
\label{rem:whatsnew}
Not the second fixed point in type $B$: that is outside the domain.  What
remains, and is genuine:
\begin{itemize}
\item \textbf{$d\ge3$ and multi-orbit $\bD$}, within the admissible domain:
  e.g.\ $\bD=\{x,x^{-1},y,y^{-1}\}$, where Motegi's $d=2$ proposition applies
  only after being iterated, and $\bD=\{1,x,x^{-1}\}$, mixing orbit types.
  Measured green.
\item \textbf{The second component $\{w_m=-1\}$ in type $C$}, which is a root
  hyperplane there and is \emph{not} $u_n=1$.
\item \textbf{Both the unification and the obstruction}: that the three families
  are the three $\iota$-orbit types (hence three and no more), and that the
  domain is cut out by $h$ in the head factor
  (Corollary~\ref{cor:domain}).  Neither is visible from two separate
  propositions.
\end{itemize}
\end{rem}

%=====================================================================
\section{The named risk: long and short exchange, and nothing happens}
\label{sec:vocab}

The brief flagged a specific danger.  The root $e_a+e_b$ is \textbf{long in
$B_N$ and short in $C_N$} (verified from $|\alpha|^{2}$ at $n=2,3,4$ on 10-18;
the naming defect I corrected in 35 places), while the boundary root is $2e_m$
(long) in $C_N$ and $e_m$ (short) in $B_N$.  If the orbit type attaches to the
\emph{root} and the multiplier to the \emph{length}, the two types must disagree
about which family gets which multiplier, and the type-$C$ theorem's vocabulary
breaks.

\textbf{It does not happen, and the proof says why.}  Reading
Theorem~\ref{thm:gaugeB}, Lemma~\ref{lem:typecompare} and
Lemma~\ref{lem:step4}: $\rho^{\alpha}$ is a function of exactly two inputs, (i)
the $\iota$-orbit decomposition of $\bD(\alpha)$, and (ii) the exponent $h$,
which appears in the head factor $(v^{h}-1)$ and in the $u$-diagonal block of
$\EX$.  \emph{No step references a root length, or $|\alpha|^{2}$, or the
long/short dichotomy.}  The exchange of lengths is therefore invisible to the
statement.

\begin{rem}[so what is $h$?]\label{rem:h}
Worth saying once precisely, because my own \texttt{2026-10-17-typeB}
Theorem~\texttt{thm:rhoB} glosses it as ``$h$ the length of the boundary root:
$2t_c$ in type $C$, $t_c$ in type $B$'', and that will mislead whoever reads it
next.  $h_X$ is the exponent in the boundary \emph{character} $w_m^{h}$ ---
equivalently (Lemma~\ref{lem:typecompare}, Corollary~\ref{cor:ratio}) the
exponent in the head factor and in $\EX$'s $u$-diagonal block.  It is not a
length: $|2e_m|^{2}=4$ and $|e_m|^{2}=1$, neither of which is $h$.  It
correlates with ``long'' within the boundary family by accident of there being
two types; and the uniform formula applies $h$ to the \emph{pair} family too,
where $e_a+e_b$ is long in $B$ and short in $C$ --- the opposite ordering.  The
formula is right; the gloss is a quantifier error of the same kind as the
long-root/short-root divisor naming, and I retract it here rather than leave it
standing in a paper Rick has.
\end{rem}

%=====================================================================
\section{What is already in Motegi, including for type $C$}
\label{sec:priority}

Task~0 of the brief was mandatory and it fired.  Reading \texttt{2610.15344}
\S``Recursive relations and stability'' for both types, first-hand:

\begin{description}
\item[Type $B$, pair family.]  \texttt{normalizedsusythree}, L2767--2771.
  \emph{Assume $r\le n-2$.  Then}
  $\widetilde{WB}_{N,n,r}(\mathbf u|\mathbf w|\bI)\big|_{u_iu_j=1}
   =\frac{(1-q)^{2}}{(1-qu_i)(1-qu_j)}
    \widetilde{WB}_{N,n-2,r}(\mathbf u\sm\{u_i,u_j\}|\mathbf w|\bI)$.
\item[Type $B$, boundary family.]  \texttt{stabilityWB}, L2779--2783.
  \emph{Assume $r\le n-1$.  Then}
  $\widetilde{WB}_{N,n,r}\big|_{u_n=1}
   =\widetilde{WB}_{N,n-1,r}(\mathbf u\sm\{u_n\})$: multiplier $1$.
\item[Type $C$, pair family.]  \texttt{normalizedsusytwo}, L2618--2622, with a
  proof: multiplier $\frac{(1-q)^{2}}{(1-qu_i^{2})(1-qu_j^{2})}$.
\item[Type $C$, boundary family.]  \texttt{stabilityWC}, L2643--2647:
  multiplier $1$ at $u_n=1$.
\end{description}

\noindent
These are $u$-specialisations with $\mathbf w$ free and independent of $\mathbf
u$; mine are $w$-specialisations with $\mathbf u=\mathbf w_{\bM}$ determined by
$\mathbf w$.  The implication runs one way, and easily:

\begin{prop}[Motegi's propositions imply the pair and boundary cases]
\label{prop:specialise}
Let $\alpha=e_a+e_b$ resp.\ $e_m$, $\bM\supseteq D(\alpha)$,
$r\le|s_\alpha\bM|$.  Then $(\mathbf u,\mathbf w)=(\mathbf w_{\bM},\mathbf w)$
with $\mathbf w$ generic in $H_\alpha$ lies on Motegi's locus $\{u_iu_j=1\}$
resp.\ $\{u_n=1\}$, with $i,j$ resp.\ $n$ the slots carrying $D(\alpha)$, and
both sides of his identity are regular there; so his identity specialises to
Corollary~\ref{cor:rhoB} for those two families, with $n=|\bM|$.
\end{prop}

\begin{proof}
On $\{w_aw_b=1\}$ the slots carrying $w_a,w_b$ satisfy $u_iu_j=1$; on
$\{w_m=1\}$ the slot carrying $w_m$ has value $1$.  Regularity: the only
denominators of \eqref{eq:T} are the $(v_p-v_{p'})$ of lines 2a and 3a,
nonvanishing because the values of $\mathbf w_{\bM}$ are pairwise distinct at
the generic point of $H_\alpha$; the $W^{B}$ prefactor is regular and nonzero
there by Lemma~\ref{lem:typecompare} (its only singularity is at $u_k=-1$,
which is not on $H_\alpha$ for a root of $B_N$); and $\EX$ has neither a zero
nor a pole there, its factors being $1-q\cdot(\text{Laurent monomial})$
evaluated away from the corresponding $(Q)$-divisor.  His hypotheses $r\le n-2$
resp.\ $r\le n-1$ are exactly $r\le|s_\alpha\bM|$.
\end{proof}

\begin{retr}[priority: $\rho^{C}$ is Motegi's, and my own paper does not say so]
\label{retr:priority}
\texttt{2026-10-18-c2} \S\texttt{sec:priority} lists under ``\textbf{Not new}''
that $\rho^{C}=(1-q)^{2}/((1-qw_a^{h})(1-qw_b^{h}))$ ``is Rick's Lemma QC and my
\texttt{2026-10-16-c2} Thm \texttt{collapse}''.  That attribution is
\textbf{incomplete and I retract it}: by Proposition~\ref{prop:specialise} in
type $C$, the same number is Motegi's \texttt{normalizedsusytwo}, L2618--2622,
proved there.  The same applies to the type-$C$ boundary multiplier
(\texttt{stabilityWC}) and, as of today, to both non-trivial type-$B$
multipliers, one of which is my own \texttt{2026-10-17} Theorem
\texttt{thm:rhoB}.

The measurement of how I missed it is unflattering and worth recording.
\texttt{stabilityWC} is L2645; my engine's \texttt{Atilde} docstring has cited
its iterate, L2822--2828, since 2026-10-15.  \texttt{normalizedsusytwo} is
\textbf{$27$ lines above the statement I was already citing}, in the same
subsection, under a heading that names it.  This is the third consecutive
session in which something I proved, or filed as a gap, sat in a paper I hold at
\texttt{deep-read} within a few lines of a locator I had already written down
(10-18: G3 at L3302--3309, four lines from Rick's citation; 10-18: the Koncki
table, reachable from a second bibitem key).  The pattern is not ``I have not
read the paper''.  It is that \textbf{an extraction level records provenance
depth, not coverage}, and the cheapest available instrument --- reading the
neighbourhood of a locator I am already using --- is the one I keep not reaching
for.

\textbf{What today's proof still buys, stated without inflation.}  A
self-contained derivation, from the symmetrisation formula, of the arbitrary-%
admissible-$\iota$-closed law in type $B$, which specialises to both of his
propositions; plus Remark~\ref{rem:whatsnew}'s three items; plus the structural
facts of \S\ref{sec:origin} and Corollary~\ref{cor:domain}, which no pair of
separate propositions exhibits.  Had I graded G2 a corollary at Task~0 and
stopped --- which the brief explicitly authorised --- I would not have found
Theorem~\ref{thm:necessary} or Corollary~\ref{cor:domain}, and I would still
believe Remark~\ref{rem:G5}'s type-$C$ answer was the general one.
\end{retr}

%=====================================================================
\section{Verification}
\label{sec:verif}

All computation is pure SymPy (no combinatorics library), in
\texttt{proofs/code-1019-typeB-gauge/}.  The engine
\texttt{q484\_engine.py} is the 2026-10-18 copy, md5
\texttt{f53c025810f95130ccec70122b1e1ae3}, identical across all five
in-repository copies (verified as a set, not by reading one).  The e-print
source used for every transcription is md5 \texttt{c15a23f2\dots}, identical
across all five local copies.

\paragraph{V3 --- the engine, validated against type-$B$ theorems not used to
write it.}  The 2026-10-17 defect was an engine that normalised type $B$ by
$\EC$ and produced a type-$C$ arm that was green and structurally blind.  The
fix is to validate against a \emph{different} theorem from the definition the
code was written from.  \textbf{Deliberately excluded} from this file:
\texttt{stabilityWB} and \texttt{normalizedsusythree}, because those two
\emph{are} (specialisations of) today's theorem, and using them to validate the
engine and then the engine to corroborate the theorem would be circular.
\begin{center}
\begin{tabular}{@{}llr@{}}
\toprule
validator & Motegi locator & result\\
\midrule
diagonal is a root product, type $B$ & \texttt{typeBmainproperties}(i), L1386 & $12/12$\\
$w$-symmetry $\widetilde{WB}(u|w|\bI)=\widetilde{WB}(u|s_iw|\bI)$
  & \texttt{prop:B\_si}(ii), L2733--2737 & $13/13$\\
boundary shift, multiplier $\frac{1-q}{1-qw_1}$
  & \texttt{prop:B\_shift}(ii), L2757--2760 & $11/11$\\
\bottomrule
\end{tabular}
\end{center}
The first of these is itself a prediction and not a transcription: I took type
$C$'s explicit diagonal and replaced \emph{only} the boundary-root factor
$(1-w_{i_a}^{-2})$ by $(1-w_{i_a}^{-1})$, and that matched.

\paragraph{Plants on V3, with admissibility decided by arithmetic.}
Type-$C$ boundary root planted into the type-$B$ diagonal: fired $10/10$.
Type-$C$ shape planted into the boundary-shift multiplier: fired $11/11$.
Third plant, \textbf{reported in full because it did not fire everywhere}: feed
the engine $W^{B}/\EC$ --- the 2026-10-17 defect itself --- and compare against
the type-$B$ diagonal.  It fired $8/10$.  The two non-firing arms are
$\bM=\bI=[N]$ at $N=2$ and $N=3$, and they are \textbf{inadmissible, not
blind}: by Remark~\ref{rem:boxed}, $\EB/\EC=\prod_{k\notin\bM}
(1-qw_k^{2})/(1-qw_k)$, an \emph{empty} product when $\bM=[N]$, so there
$\EB=\EC$ exactly and the planted object equals the true one.  Verified
symbolically, $28/28$ for the ratio and $\EB=\EC$ exactly on the $\bM=[N]$
stratum.  \textbf{So the 2026-10-17 defect is invisible on the full stratum} ---
which is a fact about where that error could not have been caught, and it is
why the plant is reported as $8/10$ with a mechanism rather than as $10/10$
after quietly dropping two arms.

\paragraph{V1 --- the headline, two mechanisms.}  $\rho^{\alpha}$ measured at
every positive root of $B_2$ and $B_3$ (nine roots: three $A$-type, three
boundary, three pair), type-$B$ normalisation, every $\bM$ in the colliding
stratum and every $\bI$, against $\prod_{m\in D(\alpha)}\psi_B(w_m)^{-1}$.
Mechanism M1 restricts to $H_\alpha$ and then divides; mechanism M2 divides as
rational functions and then restricts along a uniformiser $t$, so that a $0/0$
on $H_\alpha$ is resolved by cancellation in $t$ rather than silently.  All
three divisors admit a \emph{rational} parametrisation ($w_j=w_i$; $w_m=1$;
$w_b=1/w_a$), so no radical uniformiser is introduced --- the 2026-10-18 trap in
which one family's order came out $2$ by one script and $0$ by another.
\textbf{Result: $104/104$ discriminating instances, M1 and M2 each agreeing with
the prediction and with each other.}  Separately, $124/124$ instances of the
vacuous stratum $r>|s_\alpha\bM|$ vanish on both sides
(Proposition~\ref{prop:vac}), and $40/40$ further instances have
$\At_{s_\alpha\bM,\bI}=0$ for triangularity reasons; these $164$ are reported as
accounting, \emph{not} scored as content, since $0=\rho\cdot0$ carries no bits
about $\rho$.

\paragraph{Plants on V1.}  Three, each fired $13/13$ on its admissible domain:
$\psi$ replaced by the type-$C$ shape $(1-qw^{2})/(1-q)$; multiplier forced to
$1$; $\psi$ in place of $\psi^{-1}$.  Each of the three is
\textbf{inadmissible on nine of the thirteen root-family arms}: the three
$A$-type and six boundary arms (across $N=2,3$) have true multiplier $1$ and
$\psi_B(1)=1$, so all three plants coincide with the truth there and cannot
fire.  Only the four pair-family arms are admissible, and on them the plants
fire $13/13$.  Reported this way rather than as $13/13$ after quietly dropping
nine arms --- an inadmissible plant and a blind instrument produce the same
silence.

\paragraph{V4 --- type-independence of the unnormalised ratio
(Lemma~\ref{lem:typeindep}).}  Two mechanisms: (S) $w$ symbolic, reaching
$|\bR|+|\bD|\le3$; (Nw) $w$ numeric at a random point of the torus with $q$
symbolic, reaching $6$.  Results: $\bD=\{1\}$, $18/18$ under (S) and $18/18$
under (Nw); $\bD=\{x,x^{-1}\}$, $7/7$ under (S) and $18/18$ under (Nw);
$\bD=\{x,x^{-1},y,y^{-1}\}$, $18/18$ under (Nw) only --- $n=6$ is out of reach
symbolically, and I say so rather than imply two readings where there is one.
The plant --- prefactor \emph{not} cancelled by line 3c, i.e.\ an extra
$\prod_y(y+1)$ --- fired $18/18$ on every case including those at $-1$.  The
runs at $\bD\ni-1$ returned $0/18$, and \textbf{that reading was wrong}: see
Remark~\ref{rem:cancelfirst}.

\paragraph{V2 --- the general $\iota$-closed law, split by $r$.}  Two mechanisms
as in V4, over $N=2,3$, $|\bR|\le2$ and every $\bI$, every instance
discriminating (base $\ne0$).  Since $r=0$ is the stratum with no head slot, the
census is reported split on $r$; that split is the whole content.
\begin{center}
\begin{tabular}{@{}lcccc@{}}
\toprule
$\bD$ & admissible? & $r=0$ & $r\ge1$ & mechanisms\\
\midrule
$\varnothing$                & yes & $6/6$ & $14/14$ & (S), (Nw)\\
$\{1\}$                      & yes & $6/6$ & $14/14$ & (S), (Nw)\\
$\{x,x^{-1}\}$               & yes & $6/6$ & $14/14$ & (Nw); (S) $4/4$, $5/5$\\
$\{x,x^{-1},y,y^{-1}\}$      & yes & $6/6$ & $14/14$ & (Nw) only\\
\midrule
$\{-1\}$                     & \textbf{no} & $6/6$ & $\mathbf{0/14}$ & (S), (Nw)\\
$\{1,-1\}$                   & \textbf{no} & $6/6$ & $\mathbf{0/14}$ & (Nw); (S) $4/4$, $0/5$\\
$\{-1,x,x^{-1}\}$            & \textbf{no} & $6/6$ & $\mathbf{0/14}$ & (Nw) only\\
\bottomrule
\end{tabular}
\end{center}
\textbf{Every admissible $\bD$ is green at every $r$; every inadmissible $\bD$ is
green at $r=0$ and $0/14$ at $r\ge1$.}  The partition of the census by
Definition~\ref{def:adm} is exact, and the $r$-split inside it is exactly the
presence of a head slot.  The $d=4$ row is the case
Remark~\ref{rem:whatsnew} claims as content beyond Motegi's $d=2$ proposition
otherwise than by iteration; $d=3$ with a mixed orbit type appears in the
inadmissible block, where it also obeys the split.  Three plants --- type-$C$
shape, multiplier forced to $1$, and $\psi$ for $\psi^{-1}$ --- each fired
$20/20$ on every case; each is \textbf{inadmissible at $\bD=\{1\}$}, where the
true multiplier is $1$ and all three plants coincide with it, and that row is
excluded by arithmetic rather than scored.

\paragraph{V6 --- the cancelled form, and the refutation.}  The cancelled form
of Lemma~\ref{lem:typecompare} agrees with the engine's literal $W^{B}$ at
generic $u$, $18/18$ --- which is the check that licenses using it at $u=-1$.
Head-slot term at a $\tau$-fixed point ($N=2,3$, $r=1,2$, every $\bI$;
nine instances): vanishes $9/9$ at $+1$, fails to vanish $9/9$ at $-1$ --- a
clean split on the nose of the stated discriminator.  Adjunction law at
$\bD=\{-1\}$, symbolic in $w$ and $q$, $N=2,3$, $|\bR|\le2$, every $\bI$:
\textbf{$r=0$ gives $6/6$ and $r\ge1$ gives $\mathbf{0/14}$}
(Theorem~\ref{thm:necessary}; see also V2).  On the smaller domain
$|\bR|\le1$ the same phenomenon reads $4/9$, the four greens again being
exactly the $r=0$ instances.  Adjunction law at $\bD=\{1\}$: $9/9$; at $\bD=\{x,x^{-1}\}$: $9/9$;
at $\bD=\{1,-1\}$: $4/9$, again exactly the $r=0$ instances.  One V6 green was
initially reported as carrying no information --- $\bD=\{-1,x,x^{-1}\}$ at
$2/2$, both instances at $r=0$ because $d=3$ with $|\bR|+|\bD|\le3$ forces
$|\bR|=0$ --- and that reading is \textbf{superseded} by V2's wider
numeric-$w$ domain, which reaches $r\ge1$ there and returns $0/14$, as
Lemma~\ref{lem:head} requires.  Every green and every red falls on the side the
lemma predicts.
\end{document}
